\documentclass[11pt]{article}
\usepackage[T1]{fontenc}
\usepackage[utf8]{inputenc}
\usepackage{lmodern,amsmath,amssymb,amsthm,mathtools}
\usepackage[margin=27mm]{geometry}
\usepackage{microtype,booktabs,longtable,array,needspace}
\usepackage[colorlinks=true,linkcolor=blue,citecolor=blue,urlcolor=blue]{hyperref}
\hypersetup{pdfauthor={Tingyu Zhang},pdftitle={Sharp swap-endpoint geometry and an analytic boundary for two-state quantum embedding}}
\newtheorem{theorem}{Theorem}[section]
\newtheorem{proposition}[theorem]{Proposition}
\newtheorem{lemma}[theorem]{Lemma}
\newtheorem{corollary}[theorem]{Corollary}
\theoremstyle{definition}
\newtheorem{definition}[theorem]{Definition}
\theoremstyle{remark}
\newtheorem{remark}[theorem]{Remark}
\DeclareMathOperator{\tr}{tr}
\DeclareMathOperator{\rank}{rank}
\DeclareMathOperator{\adj}{adj}
\DeclareMathOperator{\diag}{diag}
\DeclareMathOperator{\Rea}{Re}
\DeclareMathOperator{\Ima}{Im}
\newcommand{\Q}{\mathcal Q_2}
\newcommand{\Diss}{\mathcal D}
\newcommand{\ee}{\mathbf e_z}
\newcommand{\R}{\mathbb R}
\newcommand{\C}{\mathbb C}
\newcommand{\PSD}{\succeq0}
\newcommand{\id}{\mathrm{id}}
\title{Sharp swap-endpoint geometry and an analytic boundary\newline for two-state quantum embedding}
\author{Tingyu Zhang}
\date{}
\begin{document}
\maketitle
\begin{abstract}
We study the stochastic matrices obtained by preparing and measuring a qubit in one fixed orthonormal basis, with a time-independent Gorini--Kossakowski--Sudarshan--Lindblad generator between preparation and measurement. At the deterministic swap, we determine the exact tangent cone of the closure of the realizable set. Its two nontrivial slopes are controlled by
\(r_*=(\pi-2\sqrt2)/(\pi+2\sqrt2)\).
The proof treats arbitrary Hamiltonian rotation numbers and arbitrary finite collections of jumps. We then show that the true lower boundary is the restriction of a real-analytic function near the swap. Small positive boundary points are attained by finite generators; after a measurement-preserving gauge change, their canonical noise has rank one and the unique local stationary branch lies in a common invariant Bloch plane. The argument combines a uniform weak-dissipation estimate, a spectral gap, positive-semidefinite optimality, and an analytic implicit-function theorem in the full complex rank-one chart. We also give an exact drift ellipsoid and filled-section description, a logarithmic reset construction with coefficient \(4/e\), and an explicit proof that squares of quantum channels realize every two-state stochastic matrix. The principal statements and the bridges to the original GKSL model are checked in Lean and mathlib; an explicit correspondence table records the formal scope.
\end{abstract}

\section{Introduction and setting}\label{sec:introduction}

A classical stochastic matrix can arise from a quantum evolution even when it cannot arise from a classical continuous-time Markov chain. The distinction already occurs for two outcomes. Fix the computational basis of \(\C^2\), and write a column-stochastic matrix as
\begin{equation}\label{eq:stochastic}
 T(a,b)=\begin{pmatrix}a&1-b\\1-a&b\end{pmatrix},\qquad 0\leq a,b\leq1.
\end{equation}
For a quantum channel \(\Phi\), its measured table is
\begin{equation}\label{eq:readout}
 a=\langle0|\Phi(|0\rangle\langle0|)|0\rangle,
 \qquad b=\langle1|\Phi(|1\rangle\langle1|)|1\rangle.
\end{equation}
Only the preparation and final measurement use this basis. Coherences during the evolution are unrestricted.

We consider the time-independent GKSL generators \cite{GKS1976,Lindblad1976}
\begin{equation}\label{eq:gksl}
 \mathcal L(\rho)=-i[H,\rho]+\sum_{\nu=1}^N\Diss[A_\nu](\rho),\qquad
 \Diss[A](\rho)=A\rho A^\dagger-\tfrac12\{A^\dagger A,\rho\},
\end{equation}
where \(H=H^\dagger\), \(N\) is finite, and the observation time is nonnegative. Let \(\mathcal Q_2^{\mathrm{fin}}\) be the set of tables obtained from \(e^{t\mathcal L}\), and set
\begin{equation}\label{eq:Qdefinition}
 \Q=\overline{\mathcal Q_2^{\mathrm{fin}}}\subseteq[0,1]^2.
\end{equation}
The closure includes infinite-time limits but does not change the physical model of the finite approximants. Rescaling \(H\) to \(tH\) and the jumps to \(\sqrt t A_\nu\) reduces every finite-time realization to time one. We use that normalization throughout.

The point \((0,0)\) is the deterministic swap, and \((0,1)\) is a deterministic reset. Quantum embeddability of stochastic matrices has been studied systematically in \cite{Shahbeigi2024}; the relation between classical data and coherent channel realizations is also developed in \cite{Korzekwa2018}. Our main results concern the actual boundary near the swap. Define
\begin{equation}\label{eq:fdefinition}
 f(b)=\min\{a:(a,b)\in\Q\},\qquad 0\leq b\leq1.
\end{equation}
This minimum exists: \(\Q\) is compact and contains \((b,b)\), realized by a unitary rotation. In particular \(0\leq f(b)\leq b\) and \(f(0)=0\).

Put
\begin{equation}\label{eq:constants}
 q_*=\frac{2\sqrt2}{\pi},\qquad
 r_*=\frac{1-q_*}{1+q_*},\qquad
 g=1-r_*,\qquad \beta_*=1+q_*.
\end{equation}
Below we abbreviate \(r=r_*\). We prove that the positive tangent cone of \(\Q\) at the swap is exactly
\begin{equation}\label{eq:cone-intro}
 \{(u,v)\in[0,\infty)^2:r_*v\leq u,\ r_*u\leq v\}.
\end{equation}
We also prove that there are \(b_0>0\) and a real-analytic function \(f_{\mathrm{arc}}\) on \((-b_0,b_0)\) such that
\begin{equation}\label{eq:analytic-intro}
 f(b)=f_{\mathrm{arc}}(b)\qquad(0\leq b<b_0).
\end{equation}
The existence of a single analytic stationary family is not enough for this conclusion. We show that every small boundary minimizer enters its fixed uniqueness neighborhood, after proving finite attainment and exact rank-one reduction for the canonical noise. This provides the global-to-local identification required by \eqref{eq:analytic-intro}.

Two other parts of the geometry admit explicit descriptions. At fixed linear Bloch dynamics, all admissible drifts form an ellipsoid, giving an exact interval of measured biases. Near reset, a legal one-jump family gives a relative improvement \(4/[e\log(1/\epsilon)]\), with a uniform logarithmic remainder. Finally, for the different problem of taking a square of an arbitrary quantum channel, we show directly that every table \eqref{eq:stochastic} occurs. The distinction between homogeneous semigroups and channel roots is essential \cite{WolfCirac2008}.

The proofs below are mathematical proofs in the usual notation. Section~\ref{sec:formal} identifies the corresponding Lean declarations and records the reproducible verification procedure. The analytic interval is existential; no numerical endpoint or second-order coefficient for \(f\) is asserted here.

\section{Canonical geometry and filled sections}\label{sec:geometry}

\subsection{From arbitrary jumps to a real affine equation}
Let \(X,Y,Z\) denote the Pauli matrices, and write \(\sigma=(X,Y,Z)\). The scalar part of a jump can be removed without changing the generator class. Indeed, if \(A=\alpha I+L\) with \(\tr L=0\), then
\begin{equation}\label{eq:scalar-gauge}
 \Diss[A]=\Diss[L]-i\left[\frac{i}{2}(\overline\alpha L-\alpha L^\dagger),\,\cdot\,\right].
\end{equation}
The Hamiltonian correction is Hermitian. Removing its scalar part as well, we may therefore write
\begin{equation}\label{eq:canonical}
 H=\tfrac12 h\cdot\sigma,\qquad L_\nu=v_\nu\cdot\sigma,
 \qquad K=\sum_\nu v_\nu v_\nu^\dagger\PSD,
\end{equation}
where \(h\in\R^3\) and \(K\) is a Hermitian \(3\times3\) matrix. Conversely, a positive-semidefinite \(K\) has a Gram factorization with at most three vectors, so every pair \((h,K)\) is a finite GKSL generator.

For \(\rho=(I+r\cdot\sigma)/2\), direct Pauli multiplication gives
\begin{equation}\label{eq:bloch}
 \dot r=Mr+c,\qquad
 M=[h]_\times+2(S-dI),\quad S=\Rea K,\quad d=\tr K=\tr S,
\end{equation}
where \([h]_\times r=h\times r\), and
\begin{equation}\label{eq:drift}
 c=-4(\Ima K_{23},\Ima K_{31},\Ima K_{12}),\qquad
 K=S+\frac{i}{4}[c]_\times.
\end{equation}
Indices in these formulas are \(1,2,3\). The real matrix \(S\) is positive semidefinite. The damping matrix
\begin{equation}\label{eq:damping}
 D=2(dI-S)
\end{equation}
satisfies \(D\PSD\), \(\tr D=4d\), and \(\|D\|\leq2d\). Thus the linear Bloch flow is a contraction. The full affine flow preserves the Bloch ball by complete positivity.

Set
\begin{equation}\label{eq:xz}
 x=a+b-1,\qquad z=a-b,
 \qquad a=\frac{1+x+z}{2},\quad b=\frac{1+x-z}{2}.
\end{equation}
The stochastic square becomes the diamond \(|z|\leq1-|x|\). With
\begin{equation}\label{eq:linear-readout}
 E=e^M,\qquad \ell=\int_0^1e^{tM^T}\ee\,dt,
\end{equation}
the two measured quantities are particularly simple:
\begin{equation}\label{eq:xz-readout}
 x=\ee^T E\ee,\qquad z=\ell^Tc.
\end{equation}
The first quantity depends only on \((S,h)\); the second depends linearly on the drift.

\subsection{The exact drift ellipsoid}
\begin{theorem}[Admissible drifts]\label{thm:drift-ellipsoid}
For every real symmetric \(S\PSD\),
\begin{equation}\label{eq:ellipsoid}
 S+\frac{i}{4}[c]_\times\PSD
 \quad\Longleftrightarrow\quad
 c=4\sqrt{\adj S}\,u\ \text{for some }u\in\R^3,\ \|u\|\leq1.
\end{equation}
This statement includes singular \(S\).
\end{theorem}
\begin{proof}
An orientation-preserving orthogonal change of basis diagonalizes \(S\) and transforms \([c]_\times\) covariantly. It is enough to take \(S=\diag(s_1,s_2,s_3)\), with all \(s_i\geq0\). Its two-dimensional principal minors are
\(s_i s_j-c_k^2/16\), and
\begin{equation}\label{eq:det-ellipsoid}
 \det\left(S+\frac{i}{4}[c]_\times\right)
 =s_1s_2s_3-\frac{s_1c_1^2+s_2c_2^2+s_3c_3^2}{16}.
\end{equation}
If all \(s_i>0\), nonnegativity of this determinant is equivalent to
\(\sum_i c_i^2/(16s_js_k)\leq1\), and implies every two-dimensional principal-minor inequality. This is precisely \eqref{eq:ellipsoid}. If, for example, \(s_3=0<s_1,s_2\), positivity forces \(c_1=c_2=0\) and leaves \(|c_3|\leq4\sqrt{s_1s_2}\). The adjugate is then \(\diag(0,0,s_1s_2)\), giving the same formula. If \(\rank S\leq1\), the two-dimensional principal minors force \(c=0\), while \(\adj S=0\). These cases exhaust the possibilities.
\end{proof}

\begin{proposition}[Fixed-dynamics bias interval]\label{prop:fixed-drift}
For fixed \(S\PSD\) and \(h\), the attainable values of \(z\) form exactly
\begin{equation}\label{eq:bias-interval}
 [-Z(S,h),Z(S,h)],\qquad
 Z(S,h)=4\sqrt{\ell^T(\adj S)\ell}.
\end{equation}
Every value in this interval has a finite GKSL realization. If \(Z(S,h)>0\), the upper endpoint is realized by
\begin{equation}\label{eq:optimal-drift}
 c_+=\frac{4(\adj S)\ell}{\sqrt{\ell^T(\adj S)\ell}}.
\end{equation}
\end{proposition}
\begin{proof}
By Theorem~\ref{thm:drift-ellipsoid}, the linear functional \(\ell^Tc\) is the image of the unit ball under \(u\mapsto4\langle\sqrt{\adj S}\ell,u\rangle\). Its image is the claimed interval, and its maximum is attained in the direction \(\sqrt{\adj S}\ell\). Gram factorization of the resulting \(K\) gives a finite realization.
\end{proof}

\begin{proposition}[Filling along antidiagonals]\label{prop:filling}
If \((a,b)\in\Q\) and \(-1\leq\eta\leq1\), then the point with coordinates
\(x=a+b-1\), \(z=\eta(a-b)\) also lies in \(\Q\). In particular \(\Q\) is invariant under \((a,b)\mapsto(b,a)\).
\end{proposition}
\begin{proof}
For a finite canonical generator, replace \(K\) by
\begin{equation}\label{eq:drift-scaling}
 K_\eta=\frac{1+\eta}{2}K+\frac{1-\eta}{2}\overline K.
\end{equation}
This matrix is positive semidefinite, has the same real part, and replaces \(c\) by \(\eta c\). Keep the canonical Hamiltonian fixed. Equations~\eqref{eq:xz-readout} then keep \(x\) fixed and multiply \(z\) by \(\eta\). Apply the construction to finite approximants to pass to \(\Q\). The scalar-jump correction in \eqref{eq:scalar-gauge} is made before this construction; all generators are expressed in the same measured basis.
\end{proof}

\subsection{A scalar description of the entire closed region}
\begin{theorem}[Envelope and classical half-region]\label{thm:envelope}
There is an upper-semicontinuous function \(F:[-1,1]\to[0,1]\) such that
\begin{equation}\label{eq:envelope}
 \Q=\left\{\left(\frac{1+x+z}{2},\frac{1+x-z}{2}\right):
 -1\leq x\leq1,\ |z|\leq F(x)\right\}.
\end{equation}
It satisfies \(0\leq F(x)\leq1-|x|\), \(F(-1)=0\), and
\begin{equation}\label{eq:classical-half}
 F(x)=1-x\qquad(0\leq x\leq1).
\end{equation}
Moreover it admits the relaxed variational representation
\begin{equation}\label{eq:relaxed-envelope}
 F(x)=\lim_{\delta\downarrow0}\sup\left\{
 Z(S,h):S\PSD,\ \big|\ee^Te^{[h]_\times+2(S-(\tr S)I)}\ee-x\big|<\delta
 \right\}.
\end{equation}
\end{theorem}
\begin{proof}
For every \(x\in[-1,1]\), a unitary rotation realizes \((x,z)=(x,0)\). The intersection of the compact set \(\Q\) with the line of fixed \(x\) therefore has a nonnegative maximum bias; define it to be \(F(x)\). Proposition~\ref{prop:filling} identifies this intersection with the full interval \([-F(x),F(x)]\). The diamond bound gives \(F(x)\leq1-|x|\) and \(F(-1)=0\).

If \(x_n\to x\), any subsequence along which \(F(x_n)\) converges gives, by closedness, a point \((x,\lim F(x_n))\) of the region. Its bias is at most \(F(x)\). Taking a subsequence realizing the upper limit proves upper semicontinuity.

For \(0<x\leq1\), the amplitude-damping jump \(\sqrt{-\log x}\,|0\rangle\langle1|\), with zero Hamiltonian at time one, realizes \((a,b)=(1,x)\). This has bias \(1-x\), the diamond upper bound. The limit \(x\downarrow0\) gives the remaining endpoint, proving \eqref{eq:classical-half}.

The quantities on the right of \eqref{eq:relaxed-envelope} decrease as \(\delta\downarrow0\). Each admissible \((S,h)\) has a finite endpoint realization by Proposition~\ref{prop:fixed-drift}. A sequence of nearly maximal such endpoints has a convergent subsequence in the stochastic square; the limit at fixed \(x\) has bias at most \(F(x)\). Conversely, approximate the maximal point \((x,F(x))\) by finite generators. Replacing each drift by its maximizing drift preserves its nearby \(x\) coordinate and increases its bias, so every sufficiently small neighborhood supremum is at least the approximating bias. This proves the reverse inequality and the limit formula.
\end{proof}

\section{The sharp tangent cone at the swap}\label{sec:swap-cone}

Throughout this section, \(r=r_*\) and the observation time is one. The small parameter measured at the output is \(s=a+b\). The Hamiltonian is not assumed bounded at the outset.

\subsection{Weak dissipation forced by an almost perfect swap}
\begin{lemma}[Uniform dissipative rigidity]\label{lem:rigidity}
For any finite canonical generator with output \((a,b)\), if \(s=a+b\leq1/16\), then
\begin{equation}\label{eq:rigidity}
 d\leq-\frac12\log(1-4s)\leq\frac83s.
\end{equation}
\end{lemma}
\begin{proof}
Write \(M=\Omega-D\), where \(\Omega^T=-\Omega\) and \(D=2(dI-S)\). Let
\[
 B=e^{M/2},\quad u=B\ee,\quad v=B^2\ee,\quad Q=I-B^TB.
\]
The contraction property gives \(Q\PSD\), \(\|v\|\leq\|u\|\leq1\). By the readout formula, \(\ee^Tv=-1+s\), so
\[
 \|v\|\geq1-s,\qquad \delta:=1-\|v\|^2\leq2s,
 \qquad\|v+\ee\|\leq\sqrt{2s}.
\]
Furthermore
\[
 \ee^TQ\ee+u^TQu=\delta,
 \qquad |\ee^TQu|\leq\delta/2.
\]
For \(m=\ee^Tu\), the identity
\(2m=\ee^TQu+u^T(v+\ee)\) implies
\[
 |m|\leq s/2+\sqrt{s/2}<1/4.
\]
Thus the Gram matrix \(N\) of the two columns of \(V=(\ee,u)\) satisfies \(N\succeq I/2\): its first diagonal entry is one, its second is at least \((1-s)^2\geq3/4\), and the off-diagonal entry has modulus at most \(1/4\). Consequently \(P=VN^{-1/2}\) has orthonormal columns and
\[
 \tr(P^TQP)\leq2\tr(V^TQV)=2\delta\leq4s.
\]
Let \(\sigma_1\geq\sigma_2\geq\sigma_3>0\) be the singular values of \(B\). Compression to the two-dimensional space spanned by \(P\) gives
\begin{equation}\label{eq:singular-area}
 (\sigma_1\sigma_2)^2\geq\det(P^TB^TBP)
 =\det(I-P^TQP)\geq1-4s.
\end{equation}
The last inequality uses that the two eigenvalues of \(P^TQP\) lie in \([0,1]\).

The differential inequality for the norm under \(M\) gives
\[
 \sigma_3\geq e^{-\|D\|/2}\geq e^{-\tr D/4}
 =\sqrt{\det B}.
\]
Since \(\det B=\sigma_1\sigma_2\sigma_3\), this says \(\sigma_3\geq\sigma_1\sigma_2\). Hence
\[
 e^{-2d}=\det B\geq(\sigma_1\sigma_2)^2\geq1-4s.
\]
Taking logarithms proves the first inequality in \eqref{eq:rigidity}. For \(0\leq4s\leq1/4\), integrate \((1-t)^{-1}\leq4/3\) from zero to \(4s\) to prove the second.
\end{proof}

\subsection{A comparison uniform over all Hamiltonian rotation numbers}
A rotation about the measurement axis does not change the measured table. Applying such an azimuth rotation, we may assume
\begin{equation}\label{eq:azimuth}
 h=(h_x,0,h_z),\qquad h_x\geq0.
\end{equation}
It acts orthogonally on the Bloch coordinates, transforms \(K\) by the same real congruence, preserves \(d\), and fixes \(\ee\). Define
\[
 \theta=\frac12\sqrt{h_x^2+h_z^2},\qquad
 p=\left|\langle0|e^{-iH}|0\rangle\right|^2.
\]
Both pure-Hamiltonian return probabilities equal \(p\). When \(\theta>0\), with \(n=h/(2\theta)\),
\begin{equation}\label{eq:unitary-p}
 p=\cos^2\theta+n_z^2\sin^2\theta
 =n_z^2+(1-n_z^2)\cos^2\theta.
\end{equation}
For \(j=0,1,2,\ldots\), put \(\omega_j=(2j+1)\pi\) and
\begin{equation}\label{eq:QaQb}
 C_0=\diag(1,\tfrac12,\tfrac12),\qquad
 F_j=\begin{pmatrix}0&0&2i/\omega_j\\0&0&0\\-2i/\omega_j&0&0\end{pmatrix},
 \qquad Q_{a,j}=C_0+F_j,\quad Q_{b,j}=C_0-F_j.
\end{equation}

\begin{lemma}[Uniform first-order comparison]\label{lem:uniform-swap}
If \(a+b=s\leq1/64\), then, in the gauge \eqref{eq:azimuth}, there is an integer \(j\geq0\) such that
\begin{gather}\label{eq:angle-comparison}
 p\leq\frac{35}{3}s,\qquad
 |\theta-\omega_j/2|\leq2\sqrt p,\qquad |h_z|\leq2\theta\sqrt p,\\
 \label{eq:uniform-first-order}
 a=p+\tr(Q_{a,j}K)+R_a,\qquad
 b=p+\tr(Q_{b,j}K)+R_b,\\
 \label{eq:uniform-error}
 |R_a|,|R_b|\leq34d\sqrt p+8d^2.
\end{gather}
\end{lemma}
\begin{proof}
For clarity, we give both the norm estimate and the exact integral behind the matrices in \eqref{eq:QaQb}. Encode a Hermitian matrix as \((u,r)\mapsto(uI+r\cdot\sigma)/2\). Its trace norm is \(\max(|u|,\|r\|)\). The dissipative part of the generator is
\[
 B_D(u,r)=(0,-Dr+uc).
\]
In addition to \(\|D\|\leq2d\), positivity gives \(\|c\|\leq2d\). One way to see the latter is to factor \(K=\sum v_\nu v_\nu^\dagger\), write \(v_\nu=x_\nu+iy_\nu\), and use
\(\|c\|\leq4\sum\|x_\nu\|\|y_\nu\|\leq2\sum\|v_\nu\|^2\).
Thus \(\|B_D\|\leq4d\). The Hamiltonian and full completely positive trace-preserving semigroups are contractions in this norm. Duhamel's formula therefore gives a probability error at most \(4d\) from the Hamiltonian evolution. In particular
\[
 p\leq a+4d\leq s+4d\leq35s/3.
\]
A second Duhamel expansion gives a remainder of norm at most \((4d)^2/2=8d^2\) after the first dissipative integral.

Equation~\eqref{eq:unitary-p} implies \(|n_z|\leq\sqrt p\). Choose the odd half-turn nearest to \(\theta\), equivalently choose \(j\) so that \(j\pi\leq\theta<(j+1)\pi\). If \(\delta=\theta-(j+1/2)\pi\), then \(|\delta|\leq\pi/2\) and \(|\sin\delta|=|\cos\theta|\leq\sqrt p\). The elementary bound \(|\delta|\leq(\pi/2)|\sin\delta|\) gives \(|\delta|\leq2\sqrt p\). This also rules out \(\theta=0\) in the present small-output range.

Let \(R(t)\) be the actual Hamiltonian Bloch rotation and \(R_j(t)\) the rotation about \(X\) with angular velocity \(\omega_j\). An orthogonal change of axis carrying \((1,0,0)\) to \(n=(n_x,0,n_z)\) has distance at most \(2|n_z|\) from the identity. Conjugating a rotation by this change of axis costs at most \(4|n_z|\) in operator norm. Changing the rotation angle from \(2\theta t\) to \(\omega_jt\) costs at most \(2|\theta-\omega_j/2|t\). Hence
\begin{equation}\label{eq:propagator-comparison}
 \sup_{0\leq t\leq1}\|R(t)-R_j(t)\|\leq8\sqrt p.
\end{equation}
This estimate has no factor depending on \(j\).

The first-order integral contains two propagator factors and one \(B_D\). Replacing both propagators using \eqref{eq:propagator-comparison} changes its norm by at most \(64d\sqrt p\). The resulting vector has zero trace coordinate, on which either probability readout has norm \(1/2\). The probability error is therefore at most \(32d\sqrt p\), and in particular at most \(34d\sqrt p\).

It remains to evaluate the reference first-order integral. For one traceless jump \(V=v_xX+v_yY+v_zZ\), the reference unitary is \(U_j(t)=e^{-i\omega_jtX/2}\). Since \(U_j(1)\) swaps the two basis states, the anticommutator contributions to a return probability vanish at first order. The jump amplitudes, up to constant phases, are
\[
 v_x+i(v_y\cos\omega_jt+v_z\sin\omega_jt),\qquad
 v_x-i(v_y\cos\omega_jt+v_z\sin\omega_jt).
\]
Using
\[
 \int_0^1\cos\omega_jt\,dt=0,\quad
 \int_0^1\sin\omega_jt\,dt=\frac2{\omega_j},\quad
 \int_0^1\cos^2\omega_jt\,dt=\int_0^1\sin^2\omega_jt\,dt=\frac12,
\]
with the mixed integral zero, gives
\begin{equation}\label{eq:quadratic-forms}
 \begin{aligned}
 v^\dagger Q_{a,j}v&=|v_x|^2+\frac{|v_y|^2+|v_z|^2}{2}
 -\frac4{\omega_j}\Ima(\overline v_xv_z),\\
 v^\dagger Q_{b,j}v&=|v_x|^2+\frac{|v_y|^2+|v_z|^2}{2}
 +\frac4{\omega_j}\Ima(\overline v_xv_z).
 \end{aligned}
\end{equation}
Summing over a Gram factorization of \(K\) gives the two traces in \eqref{eq:uniform-first-order}. Combining this with the second-order Duhamel bound proves \eqref{eq:uniform-error}.
\end{proof}

\subsection{The sharp spectral inequality and a matching family}
\begin{lemma}[All-frequency spectral gap]\label{lem:spectrum}
Let \(S_j=Q_{a,j}-rQ_{b,j}\), and define
\begin{equation}\label{eq:kernel}
 w_0=\frac1{\sqrt3}(-i,0,\sqrt2)^T,\qquad
 P_0=w_0w_0^\dagger,\qquad P_\perp=I-P_0.
\end{equation}
Then
\begin{equation}\label{eq:spectral-gap}
 S_0\succeq\frac g2P_\perp,\quad \ker S_0=\C w_0,
 \qquad S_j\succeq\frac g3I\quad(j\geq1).
\end{equation}
For every \(j\geq0\), both \(Q_{a,j}\succeq rQ_{b,j}\) and \(Q_{b,j}\succeq rQ_{a,j}\) hold.
\end{lemma}
\begin{proof}
The Hermitian matrix \(C_0^{-1/2}F_jC_0^{-1/2}\) has eigenvalues
\(0,\pm2\sqrt2/\omega_j\). The identity \((1+r)q_*=g\) gives
\[
 S_j=gC_0+(1+r)F_j.
\]
At \(j=0\), direct diagonalization gives the eigenvalues \(0,g/2,3g/2\), with zero eigenspace \(\C w_0\). For \(j\geq1\), the whitened perturbation has norm at most \(q_*/3\), hence \(S_j\succeq(2g/3)C_0\succeq(g/3)I\). Changing \(F_j\) to \(-F_j\) proves the reverse inequality.
\end{proof}

\begin{lemma}[Finite matching curves and exact small ordinates]\label{lem:matching}
Put
\begin{equation}\label{eq:matching-jumps}
 V_-=-\frac{i\sqrt2}{4}X+\frac12Z,\qquad V_+=V_-^\dagger,
 \qquad \alpha_0=\frac{1-q_*}{4},\quad \beta_0=\frac{1+q_*}{4}.
\end{equation}
The generator with \(H=\pi X/2\) and the single jump \(\sqrt\epsilon V_-\) has output \((a_\epsilon,b_\epsilon)\) satisfying
\begin{equation}\label{eq:matching-errors}
 |a_\epsilon-\alpha_0\epsilon|\leq2\epsilon^2,
 \qquad |b_\epsilon-\beta_0\epsilon|\leq2\epsilon^2.
\end{equation}
Every \(0<b\leq\beta_0^2/8\) is exactly attained as its second coordinate for some \(0<\epsilon\leq2b/\beta_0\), and
\begin{equation}\label{eq:matching-upper}
 f(b)<rb+40b^2.
\end{equation}
\end{lemma}
\begin{proof}
For this generator, \(d=3\epsilon/8\), and the reference Hamiltonian is exact. The quadratic forms \eqref{eq:quadratic-forms} give first-order coefficients \(\alpha_0,\beta_0\). The second-order error is at most \(8d^2=9\epsilon^2/8\leq2\epsilon^2\). Notice that \(\alpha_0=r\beta_0\).

Fix \(b\) in the stated range and let \(\epsilon_1=2b/\beta_0\). Then
\[
 b_{\epsilon_1}\geq2b-8b^2/\beta_0^2\geq b.
\]
Since \(b_0=0\) and the matrix exponential is continuous, the intermediate value theorem gives an exact \(\epsilon\in(0,\epsilon_1]\) with \(b_\epsilon=b\). Consequently
\[
 a_\epsilon-rb\leq2(1+r)\epsilon^2
 \leq\frac{8(1+r)}{\beta_0^2}b^2
 =32(1+r)^3b^2<40b^2.
\]
Here \(0<r<1/16\), as follows directly from its definition. This is a finite admissible competitor for \(f(b)\).
\end{proof}

\begin{theorem}[Quadratic outer estimate and exact swap cone]\label{thm:swap-cone}
Every \((a,b)\in\Q\) with \(a+b\leq1/256\) satisfies
\begin{equation}\label{eq:outer}
 a\geq rb-2600(a+b)^2,\qquad
 b\geq ra-2600(a+b)^2.
\end{equation}
With the convention
\[
 T^+_{\Q}(0)=\{v:\exists t_n\downarrow0,\ z_n\in\Q,
 \ z_n/t_n\longrightarrow v\},
\]
where one may equivalently use arbitrary positive \(t_n\to0\),
\begin{equation}\label{eq:exact-cone}
 T^+_{\Q}(0)=\{(u,v):u,v\geq0,\ rv\leq u,\ ru\leq v\}.
\end{equation}
\end{theorem}
\begin{proof}
For a finite output with \(s\leq1/64\), Lemmas~\ref{lem:uniform-swap} and~\ref{lem:spectrum} give
\[
 a-rb\geq gp-34(1+r)d\sqrt p-8(1+r)d^2.
\]
Completing the square in \(\sqrt p\), and then using \eqref{eq:rigidity}, yields
\[
 a-rb\geq-\left(\frac{34^2(1+r)^2}{4g}+8(1+r)\right)d^2
 \geq-2600s^2.
\]
The last numerical inequality follows already from \(0<r<1/16\). The reverse estimate is identical with \(a,b\) interchanged. A point of the closure with \(s\leq1/256\) can be approximated by finite points lying eventually in the larger window \(s<1/64\); continuity proves \eqref{eq:outer}.

For a tangent sequence, divide \eqref{eq:outer} by \(t_n\). Since \(z_n/t_n\) converges, the quadratic error divided by \(t_n\) tends to zero. Positivity of the coordinates gives the outer inclusion in \eqref{eq:exact-cone}.

For the reverse inclusion, let \((u,v)\) satisfy the four displayed inequalities, and define
\[
 \lambda=\frac{\beta_0v-\alpha_0u}{\beta_0^2-\alpha_0^2},\qquad
 \nu=\frac{\beta_0u-\alpha_0v}{\beta_0^2-\alpha_0^2}.
\]
Both coefficients are nonnegative. Use \(H=\pi X/2\) and the two jumps
\(\sqrt{\epsilon\lambda}V_-\), \(\sqrt{\epsilon\nu}V_+\). Linearity of the first dissipative term and the same Duhamel bound give output
\[
 \epsilon\{\lambda(\alpha_0,\beta_0)+\nu(\beta_0,\alpha_0)\}+O(\epsilon^2)
 =\epsilon(u,v)+O(\epsilon^2).
\]
These are actual finite GKSL points for every \(\epsilon\geq0\), which proves the inner inclusion, including the zero vector.
\end{proof}

\section{Finite minimizers and the analytic boundary}\label{sec:analytic}

The spectral equality direction in Lemma~\ref{lem:spectrum} predicts a rank-one stationary family. To identify it with the true boundary, we first localize every sufficiently good generator, including generators whose parameters initially have no bounds.

\subsection{Uniform localization of all near minimizers}
For a fixed \(C>0\), consider finite generators satisfying
\begin{equation}\label{eq:near-class}
 b>0,\qquad a\leq rb+Cb^2.
\end{equation}
The constants below depend on \(C\), but not on the generator. Define
\begin{equation}\label{eq:localization-constants}
 \begin{aligned}
 D_*&=16/3,\\
 L_C&=C+\left(\frac{34^2(1+r)^2}{2g}+8(1+r)\right)D_*^2,\\
 P_C&=2L_C/g,\qquad E_C=34D_*\sqrt{P_C}+8D_*^2.
 \end{aligned}
\end{equation}

\begin{lemma}[Uniform localization]\label{lem:localization}
For all sufficiently small \(b>0\), every generator in \eqref{eq:near-class}, after the azimuth gauge, has \(j=0\) in Lemma~\ref{lem:uniform-swap}, and
\begin{equation}\label{eq:localized}
 p\leq P_Cb^2,\qquad b/4\leq d\leq D_*b,\qquad
 \tr(P_\perp K)\leq P_Cb^2.
\end{equation}
Uniformly over all such generators as \(b\downarrow0\),
\begin{equation}\label{eq:uniform-limits}
 h\longrightarrow(\pi,0,0),\qquad
 K\longrightarrow0,\qquad
 \frac Kd\longrightarrow P_0,\qquad
 \frac db\longrightarrow\frac{1}{(2/3)\beta_*}.
\end{equation}
\end{lemma}
\begin{proof}
Decrease the allowed range of \(b\) so that \(s=a+b\leq2b\leq1/64\). Lemma~\ref{lem:rigidity} gives \(d\leq D_*b\). Keep the nonnegative spectral term in the proof of Theorem~\ref{thm:swap-cone} and complete only half of the quadratic term in \(\sqrt p\). This gives
\begin{equation}\label{eq:local-energy}
 \frac g2p+\tr(S_jK)\leq L_Cb^2.
\end{equation}
In particular \(p\leq P_Cb^2\), and the remainders in \eqref{eq:uniform-error} are at most \(E_Cb^2\). Since \(Q_{b,j}\preceq2I\),
\[
 b\leq2d+(P_C+E_C)b^2.
\]
Thus \(d\geq b/4\) once \((P_C+E_C)b\leq1/2\). If \(j\geq1\), the all-frequency gap gives
\[
 gb/12\leq gd/3\leq\tr(S_jK)\leq L_Cb^2,
\]
which is impossible for sufficiently small positive \(b\). Hence \(j=0\). Now \(S_0\succeq(g/2)P_\perp\) and \eqref{eq:local-energy} prove the remaining estimate in \eqref{eq:localized}.

The angle and axis bounds in \eqref{eq:angle-comparison} imply \(h=(\pi,0,0)+O(b)\) uniformly. To control the noise direction, write \(K\) in blocks relative to \(\C w_0\oplus w_0^\perp\). If \(t=\tr(P_\perp K)\), positivity bounds the norm of the off-diagonal block by \(\sqrt{dt}\). The perpendicular block has norm at most \(t\), and its omission changes the parallel coefficient by \(t\). Therefore
\[
 \left\|\frac Kd-P_0\right\|\leq\frac{2t}{d}+2\sqrt{\frac td}\longrightarrow0
\]
by \eqref{eq:localized}. Finally divide the second equation in \eqref{eq:uniform-first-order} by \(d\). The terms \(p/d\) and \(R_b/d\) tend to zero uniformly, whereas
\[
 \tr\left(Q_{b,0}\frac Kd\right)\longrightarrow
 \tr(Q_{b,0}P_0)=\frac23\beta_*>0.
\]
This proves the last limit in \eqref{eq:uniform-limits}.
\end{proof}

\begin{lemma}[Finite attainment near the swap]\label{lem:attainment}
For all sufficiently small \(b>0\), the point \((f(b),b)\) is attained by a finite canonical generator. It can be realized using at most three traceless jumps. Every canonical minimizer satisfies the uniform localization conclusions of Lemma~\ref{lem:localization}.
\end{lemma}
\begin{proof}
The point \((f(b),b)\) belongs to \(\Q\) by compactness. Choose finite outputs \((a_n,b_n)\) converging to it. By Lemma~\ref{lem:matching},
\(f(b)<rb+40b^2\). For fixed positive \(b\), continuity therefore implies
\[
 b_n>0,\qquad a_n\leq rb_n+41b_n^2
\]
for all sufficiently large \(n\). After applying the azimuth gauge to each approximant, Lemma~\ref{lem:localization} puts all its Hamiltonians and noise matrices in one bounded set. For example, after making the small-\(b\) interval smaller, one may use \(\|h_n\|\leq6\) and \(\tr K_n\leq1\). A subsequence converges to \((h,K)\) with \(K\PSD\). The finite-dimensional matrix exponential and the readout are continuous, so the limit generator attains \((f(b),b)\). A Gram factorization gives at most three jumps. Each minimizer itself satisfies \eqref{eq:near-class} with \(C=40\), or with \(C=41\), so the uniform conclusions apply to all of them.
\end{proof}

\subsection{Positive-semidefinite optimality forces exact rank one}
View the two canonical probabilities as real-analytic functions \(a(h,K)\), \(b(h,K)\) on the ambient real vector space of \(h\) and Hermitian \(K\). They are defined there by the same matrix exponential, even when \(K\) is not positive semidefinite. Their matrix gradients use the real inner product \(\langle A,B\rangle=\tr(AB)\).

\begin{lemma}[Exact rank-one reduction]\label{lem:rank-one}
Use the azimuth normalization \eqref{eq:azimuth}. For every sufficiently small positive \(b\), every finite canonical minimizer of \(f(b)\) has \(\rank K=1\). Its radial constraint derivative \(D_Kb[K]\) is positive, and the multiplier
\begin{equation}\label{eq:multiplier}
 \mu=\frac{D_Ka[K]}{D_Kb[K]}
\end{equation}
converges uniformly to \(r\) as \(b\downarrow0\). The dual matrix
\begin{equation}\label{eq:dual}
 \mathcal S=\nabla_Ka-\mu\nabla_Kb
\end{equation}
obeys \(\mathcal S\PSD\), \(\mathcal SK=0\), and \(\mathcal S\to S_0\).
\end{lemma}
\begin{proof}
At \(h_0=(\pi,0,0)\), \(K=0\), the exact first-order computation gives
\[
 \nabla_Ka(h_0,0)=Q_{a,0},\qquad
 \nabla_Kb(h_0,0)=Q_{b,0}.
\]
The identities hold for the entire Hermitian gradient: the jump computation holds for every complex vector, and rank-one Hermitian matrices span the real space of Hermitian matrices. Lemma~\ref{lem:localization} and continuity of the derivatives imply, uniformly over all minimizers,
\[
 D_Kb[K/d]\longrightarrow\tr(Q_{b,0}P_0)=\tfrac23\beta_*>0,
 \qquad D_Ka[K/d]\longrightarrow r\tfrac23\beta_*.
\]
Since \(d>0\), the radial derivative is positive for all sufficiently small \(b\), and \eqref{eq:multiplier} is well-defined with \(\mu\to r\). Continuity then gives \(\mathcal S\to S_0\).

We give the cone optimality argument explicitly because a rank constraint is not imposed in the minimization. Let \(W\PSD\) be arbitrary. The path \(K+t(W-K)\) remains positive semidefinite for \(0\leq t\leq1\). By the nonzero radial derivative, a real scalar factor \(\lambda(t)\), with \(\lambda(0)=1\), can be chosen by the implicit-function theorem so that
\[
 b\bigl(h,\lambda(t)(K+t(W-K))\bigr)=b(h,K).
\]
For small nonnegative \(t\), this remains an admissible generator. Minimality and differentiation give
\(\tr(\mathcal S(W-K))\geq0\).
By the definition of \(\mu\), \(\tr(\mathcal SK)=0\), so \(\tr(\mathcal SW)\geq0\) for every \(W\PSD\). Thus \(\mathcal S\PSD\). Two positive-semidefinite matrices with zero trace product satisfy \(\mathcal SK=0\), as follows by applying the zero trace to \(K^{1/2}\mathcal SK^{1/2}\).

The two positive eigenvalues of \(S_0\) are \(g/2\) and \(3g/2\). Consequently every sufficiently close \(\mathcal S\) has at least two positive eigenvalues. Its kernel has dimension at most one. Since \(K\ne0\) and its range lies in this kernel, \(\rank K=1\). The same radial correction applied to an arbitrary two-sided Hamiltonian variation proves stationarity of \(a-\mu b\) in the Hamiltonian coordinates.
\end{proof}

\subsection{The full complex rank-one analytic system}
The preceding rank-one statement concerns the canonical \(3\times3\) noise matrix. Near \(P_0\), its \(zz\) entry is positive. Every such rank-one matrix has a unique representation
\begin{equation}\label{eq:rank-one-chart}
 K=kww^\dagger,\qquad k=K_{zz}>0,\qquad
 w=(\xi+iv,q+i\zeta,-1)^T.
\end{equation}
Indeed \(w_x=-K_{xz}/K_{zz}\) and \(w_y=-K_{yz}/K_{zz}\), and the remaining entries follow from rank one. After the azimuth gauge write
\[
 H=\frac\omega2X+hZ,
 \qquad y=(\omega,h,\mu,\xi,v,q,\zeta).
\]
Lemma~\ref{lem:localization} and Lemma~\ref{lem:rank-one} give
\begin{equation}\label{eq:chart-base}
 k\longrightarrow0,\quad
 y\longrightarrow y_0=(\pi,0,r,0,1/\sqrt2,0,0),\quad
 k/d\longrightarrow2/3.
\end{equation}
This chart includes all complex jump directions near the limiting noise; no plane condition has been imposed.

For later clarity, the entire analytic system can be written as a real \(4\times4\) matrix exponential. With coordinates \((u,r_x,r_y,r_z)\), put
\begin{equation}\label{eq:augmented-Q}
 Q(k,y)=
 \begin{pmatrix}0&0&0&0\\0&0&-2h&0\\0&2h&0&-\omega\\0&0&\omega&0\end{pmatrix}
 +k\begin{pmatrix}
 0&0&0&0\\
 4\zeta&-2(W+1)&2C&-2\xi\\
 -4v&2C&-2(U+1)&-2q\\
 4(\xi\zeta-vq)&-2\xi&-2q&-2(U+W)
 \end{pmatrix},
\end{equation}
where \(U=\xi^2+v^2\), \(W=q^2+\zeta^2\), and \(C=\xi q+v\zeta\). If \(E=e^{Q(k,y)}\), indexed from zero to three, define
\begin{equation}\label{eq:analytic-probabilities}
 A(k,y)=\frac{1+E_{33}+E_{30}}2,\qquad
 B(k,y)=\frac{1+E_{33}-E_{30}}2,\qquad J=A-\mu B.
\end{equation}
These are real-analytic for all real arguments. For \(k\geq0\), they are exactly the probabilities of the Hamiltonian above and the traceless jump \(\sqrt{k}\,w\cdot\sigma\).

The stationary equations need a rescaling at \(k=0\). Set
\begin{equation}\label{eq:stationary-system}
 G(k,y)=\left(J_\omega,J_h,J_k,
 \frac{J_\xi}{k},\frac{J_v}{k},\frac{J_q}{k},\frac{J_\zeta}{k}\right).
\end{equation}
The third equation is the derivative with respect to the dissipative scale \(k\); no derivative with respect to the multiplier is required. At \(k=0\), the probabilities are independent of the four jump-direction variables. Thus every numerator in the last four entries vanishes identically on \(k=0\), and the analytic extension is
\begin{equation}\label{eq:analytic-division}
 \frac{J_u(k,y)}k=\int_0^1J_{ku}(tk,y)\,dt,
 \qquad u\in\{\xi,v,q,\zeta\}.
\end{equation}
This gives a real-analytic \(G\) on a full neighborhood of \((0,y_0)\), including negative \(k\).

\begin{lemma}[A unique analytic stationary branch]\label{lem:analytic-system}
There are open neighborhoods \(U_0\) of \(0\) and \(V_0\) of \(y_0\), and a real-analytic map \(\psi:U_0\to V_0\), such that \(\psi(0)=y_0\) and
\begin{equation}\label{eq:fixed-uniqueness}
 G(k,y)=0\quad\Longleftrightarrow\quad y=\psi(k)
 \qquad(k\in U_0,\ y\in V_0).
\end{equation}
The function \(\mathcal B(k)=B(k,\psi(k))\) satisfies
\begin{equation}\label{eq:beta-derivative}
 \mathcal B(0)=0,\qquad \mathcal B'(0)=\beta_*>0,
\end{equation}
and has a real-analytic local inverse \(\kappa\) with \(\kappa(0)=0\).
\end{lemma}
\begin{proof}
At zero dissipation the probability is
\[
 p(\omega,h)=\cos^2 R+\frac{h^2}{R^2}\sin^2 R,
 \qquad R=\sqrt{(\omega/2)^2+h^2}.
\]
Its gradient vanishes at \((\pi,0)\) and its Hessian there is
\(\diag(1/2,8/\pi^2)\). At \(\omega=\pi,h=0\), the first derivatives in \(k\), obtained from \eqref{eq:quadratic-forms}, are
\[
 A_k=C_1-4v/\pi,\qquad B_k=C_1+4v/\pi,
 \qquad C_1=\xi^2+v^2+(q^2+\zeta^2+1)/2.
\]
At \(\mu=r\) their difference is the exact square
\begin{equation}\label{eq:leading-square}
 A_k-rB_k=g\left(\xi^2+(v-1/\sqrt2)^2+\frac{q^2+\zeta^2}{2}\right).
\end{equation}
It follows that \(G(0,y_0)=0\). Order the columns of \(D_yG\) as the two Hamiltonian variables, the multiplier, and the four direction variables. The same grouping of rows gives the block lower-triangular matrix
\begin{equation}\label{eq:Jacobian}
 D_yG(0,y_0)=
 \begin{pmatrix}
 g\diag(1/2,8/\pi^2)&0&0\\
 *&-\beta_*&0\\
 *&*&g\diag(2,2,1,1)
 \end{pmatrix}.
\end{equation}
The zero blocks follow respectively from independence of the jump direction at \(k=0\), the zero unitary gradient, and the stationary point of \eqref{eq:leading-square}. In particular
\[
 \det D_yG(0,y_0)=-\frac{16\beta_*g^6}{\pi^2}\ne0.
\]
The real-analytic implicit-function theorem proves \eqref{eq:fixed-uniqueness} on fixed neighborhoods. At the base point every first derivative of \(B\) with respect to \(y\) vanishes, while \(B_k(0,y_0)=\beta_*\). The chain rule therefore gives \eqref{eq:beta-derivative}. The one-dimensional analytic inverse-function theorem supplies \(\kappa\), with both identities
\begin{equation}\label{eq:both-inverses}
 \mathcal B(\kappa(b))=b,\qquad \kappa(\mathcal B(k))=k
\end{equation}
for sufficiently small real arguments.
\end{proof}

\subsection{Identification with the true lower boundary}
\begin{theorem}[A true analytic boundary arc]\label{thm:analytic-boundary}
There exist \(b_0>0\) and a real-analytic function \(f_{\mathrm{arc}}\) on \((-b_0,b_0)\) such that
\begin{equation}\label{eq:actual-analytic-arc}
 f_{\mathrm{arc}}(b)=f(b)\qquad(0\leq b<b_0).
\end{equation}
For every sufficiently small positive \(b\), the boundary point has a finite canonical realization with noise rank one, and all azimuth-normalized minimizing generators lie on the unique branch of Lemma~\ref{lem:analytic-system}.
\end{theorem}
\begin{proof}
Choose the fixed neighborhoods and inverse from Lemma~\ref{lem:analytic-system}. For every sufficiently small positive \(b\), Lemma~\ref{lem:attainment} supplies a finite minimizer. Lemmas~\ref{lem:localization} and~\ref{lem:rank-one} apply to every such minimizer and give the uniform limits \eqref{eq:chart-base}. Hence, after making the allowed interval smaller, the scale and chart of every minimizer satisfy \(k\in U_0\), \(y\in V_0\), \(k>0\).

For completeness, the order of quantifiers here is important. If there were no common interval with this property, one could choose \(b_n\downarrow0\) and one offending minimizer at each \(b_n\). The uniform estimates \eqref{eq:localized}, convergence of the normalized radial derivatives, and the explicit chart \eqref{eq:rank-one-chart} would force its parameters to \((0,y_0)\), contradicting failure of membership in the fixed open neighborhoods. No generator is held fixed in this argument.

Optimality gives the equations \(G(k,y)=0\). The Hamiltonian equations follow from the radial constraint correction used in Lemma~\ref{lem:rank-one}. Variations of \(k>0\) and of the four complex chart coordinates give the remaining equations; the multiplier is exactly \eqref{eq:multiplier}. Thus fixed-neighborhood uniqueness gives \(y=\psi(k)\). Since the second coordinate is \(b\),
\[
 b=\mathcal B(k),\qquad k=\kappa(b)
\]
by the scalar left inverse in \eqref{eq:both-inverses}. It follows that every minimizer has
\[
 f(b)=A(\kappa(b),\psi(\kappa(b))).
\]
Define the right-hand side for all sufficiently small real \(b\). It is real-analytic by composition. At \(b=0\), \(\kappa(0)=0\), \(\psi(0)=y_0\), and \(A(0,y_0)=0=f(0)\). Shrinking a symmetric open interval around zero gives \eqref{eq:actual-analytic-arc} including its endpoint.
\end{proof}

\begin{proposition}[Exact planar symmetry of the branch]\label{prop:planarity}
For every sufficiently small positive \(b\), the canonical realization on the branch in Theorem~\ref{thm:analytic-boundary} has
\begin{equation}\label{eq:planar-branch}
 H=\frac\omega2X,\qquad
 V=\sqrt{k}\,(ivX+qY-Z),\qquad k>0.
\end{equation}
In Bloch coordinates, the plane \(r_x=0\), which contains the measurement axis, is invariant under the affine dynamics.
\end{proposition}
\begin{proof}
The transformation on the physical matrices
\[
 H\longmapsto-Z\overline HZ,\qquad V\longmapsto Z\overline VZ
\]
intertwines the evolution by \(\rho\mapsto Z\overline\rho Z\). It fixes both preparation projectors and both measurement probabilities. In the chart it is the linear involution
\[
 \mathcal R(\omega,h,\mu,\xi,v,q,\zeta)
 =(\omega,-h,\mu,-\xi,v,q,-\zeta).
\]
Thus \(A(k,\mathcal Ry)=A(k,y)\), \(B(k,\mathcal Ry)=B(k,y)\), and the stationary system transforms by the corresponding signs. The extension at \(k=0\) has the same covariance, either by \eqref{eq:analytic-division} or by continuity.

The base point is fixed by \(\mathcal R\). Intersect the uniqueness neighborhood with its reflected image and shrink the scale neighborhood so that both \(\psi(k)\) and \(\mathcal R\psi(k)\) remain in it. Both solve the same stationary system. Uniqueness forces \(\mathcal R\psi(k)=\psi(k)\), so \(h=\xi=\zeta=0\) exactly. This is \eqref{eq:planar-branch}.

Substitution in \eqref{eq:augmented-Q} gives zero drift in the \(x\) coordinate, zero coupling between \(x\) and \(y,z\), and
\(\dot r_x=-2k(q^2+1)r_x\). Therefore \(r_x=0\) is an invariant plane for the affine equation. The rank-one statement is \(\rank K=1\); it does not assert that the traceless \(2\times2\) matrix \(V\) has rank one.
\end{proof}





\section{A logarithmic reset construction}\label{sec:reset}

The reset corner has a different asymptotic scale. We give a finite construction for every sufficiently small prescribed \(\epsilon=1-b\), rather than just a sequence approaching the corner. All jumps in this section are the original \(2\times2\) matrices used in \eqref{eq:gksl}; no scalar part is removed.

For \(0<\kappa<1\), put
\begin{equation}\label{eq:reset-jump}
 m=\sqrt{\frac{1-\kappa}{2}},\qquad n=\sqrt{\frac{1+\kappa}{2}},\qquad
 J_0=\begin{pmatrix}-m&0\\n&0\end{pmatrix},\qquad
 U_r=\begin{pmatrix}\cos r&\sin r\\-\sin r&\cos r\end{pmatrix}.
\end{equation}
We use \(H=0\) and \(J=U_rJ_0U_r^T\). This is a nonzero rank-one jump. Its pure stationary state is the projector onto \(U_r|1\rangle\), because \(J_0|1\rangle=0\).

\begin{theorem}[Uniform reset construction]\label{thm:reset}
There exist constants \(C>0\) and \(\epsilon_0>0\) such that every \(0<\epsilon<\epsilon_0\) has a time-one realization \((a,1-\epsilon)\) with zero Hamiltonian and a single rank-one jump, where
\begin{equation}\label{eq:reset-remainder}
 0<a<\epsilon,\qquad
 \left|\frac a\epsilon-\left(1-\frac{4/e}{\log(1/\epsilon)}\right)\right|
 \leq C\frac{\log\log(1/\epsilon)}{\log^2(1/\epsilon)}.
\end{equation}
The jump can be chosen to be a positive real scalar multiple of \(U_rJ_0U_r^T\) from \eqref{eq:reset-jump}.
\end{theorem}
\begin{proof}
We first solve this generator exactly. Write \(P_0=|0\rangle\langle0|\), \(P_1=|1\rangle\langle1|\). The dissipator of \(J_0\) has the following three Hermitian modes:
\begin{equation}\label{eq:reset-modes}
 \begin{array}{c|c}
 \text{mode}&\text{eigenvalue}\\ \hline
 P_1&0\\
 X&-1/2\\
 \kappa(P_0-P_1)+2mnX&-(1+\kappa)/2.
 \end{array}
\end{equation}
These identities follow by multiplying the two triangular matrices in \eqref{eq:reset-jump}. Conjugation by \(U_r\) gives the modes of \(J\). Set
\begin{equation}\label{eq:reset-coefficients}
 \sigma=\sqrt{1-\kappa^2}=2mn,\quad
 B=\cos2r+\frac\sigma\kappa\sin2r,\quad A=\cos2r\,B,
 \quad S=e^{-t/2},\quad F=S^{1+\kappa}.
\end{equation}
Expanding the two input projectors in the modes \eqref{eq:reset-modes} gives the exact output
\begin{equation}\label{eq:reset-output}
 \epsilon=\sin^2r\,[1-(1+B)S+BF],\qquad
 a=\epsilon+AF-(A-1)S,\qquad b=1-\epsilon.
\end{equation}
For example, the coherence mode decays by \(S\), the fast mode by \(F\), and the stationary projector remains fixed; substituting their diagonal entries gives the two terms in brackets. At \(t=0\), these formulas give \((a,b)=(1,1)\), as required.

Suppose \(A>1\). The gain relative to the classical boundary \(a=\epsilon\) is \((A-1)S-AS^{1+\kappa}\). Its maximum as a function of \(S\in(0,1)\) occurs at
\begin{equation}\label{eq:optimal-reset-scale}
 S_* =\left(\frac{A-1}{(1+\kappa)A}\right)^{1/\kappa},
 \qquad t_*=-2\log S_*>0.
\end{equation}
At this scale the gain is
\begin{equation}\label{eq:reset-gain}
 \Delta=\epsilon-a=\kappa A S_*^{1+\kappa}>0.
\end{equation}
We now choose a curve on which this gain has the desired logarithmic size. For \(L>2\), let
\begin{equation}\label{eq:reset-schedule}
 r=e^{-L/2},\qquad \kappa=1-\frac2L,
 \qquad \sigma=\frac{2\sqrt{L-1}}L,
\end{equation}
and let \(A,B,S_*,t_*,\epsilon,\Delta\) be the preceding quantities evaluated on this curve. All estimates that follow are on a single tail \(L\geq L_0\), with constants independent of \(L\).

The identity
\begin{equation}\label{eq:reset-correction}
 A-1=\frac{2\sigma r}{\kappa}\,D_L,\qquad
 D_L=\frac{\sin2r}{2r}
 \left(\cos2r-\frac\kappa\sigma\sin2r\right)
\end{equation}
shows that \(D_L=1+O(\sqrt L\,e^{-L/2})\). In particular \(A>1\) eventually. Equations~\eqref{eq:optimal-reset-scale} and~\eqref{eq:reset-gain} give the exact factorization
\begin{equation}\label{eq:reset-prefactor}
 \frac{\Delta}{r^2}
 =\kappa\left(\frac{2\sigma}{\kappa(1+\kappa)}\right)^{1+1/\kappa}
 e^{-1/\kappa}\,
 A\left(\frac{D_L}{A}\right)^{1+1/\kappa}.
\end{equation}
The exponential factor is exact: \(r^{1/\kappa-1}=e^{-1/\kappa}\) under \eqref{eq:reset-schedule}.

Here \(A=1+O(e^{-L/2}/\sqrt L)\), while the last factor in \eqref{eq:reset-prefactor} is \(1+O(\sqrt L\,e^{-L/2})\). The logarithm of the remaining prefactor satisfies
\begin{equation}\label{eq:reset-log-expansion}
 \begin{aligned}
 &\log\kappa+\left(1+\frac1\kappa\right)
   \log\frac{2\sigma}{\kappa(1+\kappa)}-\frac1\kappa\\
 &\hspace{18mm}=\log4-\log L-1+O\left(\frac{\log L}{L}\right).
 \end{aligned}
\end{equation}
Indeed \(1/\kappa=1+O(1/L)\), \(\log\kappa=O(1/L)\), and
\(\log[2\sigma/(\kappa(1+\kappa))]=\log2-\tfrac12\log L+O(1/L)\).
The estimates follow directly from the mean value theorem for \(\log\) on a fixed neighborhood of one, so the remainders are uniform on the chosen tail. Exponentiating \eqref{eq:reset-log-expansion} in \eqref{eq:reset-prefactor} yields
\begin{equation}\label{eq:reset-gain-asymptotic}
 \frac{\Delta}{r^2}=\frac{4}{eL}+O\left(\frac{\log L}{L^2}\right).
\end{equation}
Moreover \(S_*\leq C_1r/\sqrt L\) for large \(L\): its base in \eqref{eq:optimal-reset-scale} has that bound, lies in \((0,1)\), and its exponent is greater than one. Equation~\eqref{eq:reset-output} then gives
\begin{equation}\label{eq:epsilon-equivalence}
 \frac\epsilon{r^2}=1+O(r/\sqrt L+r^2),\qquad
 E:=\log(1/\epsilon)=L+O(r/\sqrt L+r^2).
\end{equation}
Combining \eqref{eq:reset-gain-asymptotic} and~\eqref{eq:epsilon-equivalence}, and using \(a=\epsilon-\Delta\), proves
\[
 \frac a\epsilon=1-\frac4{eE}+O\left(\frac{\log E}{E^2}\right).
\]
Also \(0<\Delta<\epsilon\) on a sufficiently large tail, so \(0<a<\epsilon\).

It remains to obtain every small prescribed \(\epsilon\). The function \(\epsilon(L)\) is continuous, positive on a tail, and tends to zero by \eqref{eq:epsilon-equivalence}. Fix a point \(L_0\) beyond all the preceding thresholds. For any \(0<\epsilon<\epsilon(L_0)\), choose \(L_1>L_0\) with \(\epsilon(L_1)<\epsilon\). The intermediate value theorem on \([L_0,L_1]\) gives a parameter with exactly the desired \(\epsilon\). The uniform tail estimates apply at every such choice. Finally absorb time into the jump:
\begin{equation}\label{eq:reset-actual-unit-jump}
 J_{\mathrm{unit}}(L)=\sqrt{t_*(L)}\,
 U_{r(L)}\begin{pmatrix}-m(L)&0\\n(L)&0\end{pmatrix}U_{r(L)}^T.
\end{equation}
It remains nonzero and rank one, and its time-one output is the exact point from \eqref{eq:reset-output}. This proves all assertions.
\end{proof}

\section{Square roots of channels fill the stochastic square}\label{sec:roots}

The previous sections concern homogeneous semigroups. We now consider the measured action of a power of one arbitrary quantum channel. If \(\pi(\Phi)\) denotes the pair in \eqref{eq:readout}, define
\begin{equation}\label{eq:root-layer}
 \mathcal R_n=\{\pi(\Phi^n):\Phi\text{ is a qubit completely positive trace-preserving map}\}.
\end{equation}
The same channel is used at every step. Positivity and trace preservation give \(\mathcal R_n\subseteq[0,1]^2\).

\begin{theorem}[The square-root region]\label{thm:square-roots}
\begin{equation}\label{eq:R2-square}
 \mathcal R_2=[0,1]^2.
\end{equation}
\end{theorem}
\begin{proof}
For \(0\leq r\leq1\), define a channel \(\Phi_r\) by the Kraus matrices
\begin{equation}\label{eq:root-kraus}
 K_0=\frac1{\sqrt{1+r^2}}
 \begin{pmatrix}r&-r^2\\1&r^3\end{pmatrix},\qquad
 K_1=\begin{pmatrix}0&0\\0&\sqrt{1-r^4}\end{pmatrix}.
\end{equation}
Indeed \(K_0^\dagger K_0=\diag(1,r^4)\) and \(K_1^\dagger K_1=\diag(0,1-r^4)\), so their sum is the identity. Put
\(\psi_r=(r,1)^T/\sqrt{1+r^2}\). We have
\[
 \Phi_r(P_0)=|\psi_r\rangle\langle\psi_r|,\qquad
 K_0\psi_r=r|1\rangle,\qquad K_1\psi_r=\sqrt{1-r^2}|1\rangle.
\]
Thus \(\Phi_r^2(P_0)=P_1\). Direct substitution for the other input gives
\begin{equation}\label{eq:root-readouts}
 \pi(\Phi_r)=\left(\frac{r^2}{1+r^2},1-\frac{r^4}{1+r^2}\right),\qquad
 \pi(\Phi_r^2)=(0,1-r^4).
\end{equation}

Mix the root channel with the identity, defining
\(\Psi_{q,r}=q\id+(1-q)\Phi_r\) for \(0\leq q\leq1\). This is a channel. Its square is
\(q^2\id+2q(1-q)\Phi_r+(1-q)^2\Phi_r^2\), so its measured coordinates are
\begin{equation}\label{eq:root-square-formula}
 \begin{aligned}
 a(q,r)&=q^2+2q(1-q)\frac{r^2}{1+r^2},\\
 b(q,r)&=q^2+2q(1-q)\left(1-\frac{r^4}{1+r^2}\right)
            +(1-q)^2(1-r^4).
 \end{aligned}
\end{equation}
Fix \(a\in[0,1]\). For each \(r\), the first expression is a continuous strictly increasing function of \(q\in[0,1]\), taking the values zero and one at the endpoints. To check strict increase, write it as \((1-2t)q^2+2tq\), where \(t=r^2/(1+r^2)\in[0,1/2]\). Thus there is a unique \(q=q_a(r)\) giving the prescribed \(a\). This choice is continuous: any convergent sequence of \(r\)'s has subsequential limits of \(q\) in the compact interval, and uniqueness identifies every such limit with the value at the limiting \(r\).

At \(r=0\), \(q_a(0)=\sqrt a\) and \(b=1\). At \(r=1\), \(q_a(1)=a\) and \(b=a\). Continuity and the intermediate value theorem therefore give every \(b\in[a,1]\). This fills the upper triangle of the stochastic square. Conjugating the root channel by the basis-swap unitary interchanges \(a,b\), so its square fills the lower triangle as well. Every point is the square of a single explicitly constructed CPTP channel, proving \eqref{eq:R2-square}.
\end{proof}


\section{Formal verification and scope}\label{sec:formal}

\subsection{Definitions and semantic bridges}
The accompanying Lean development uses Lean and mathlib version \texttt{v4.35.0}. The starting object is the finite GKSL superoperator \eqref{eq:gksl}, acting on complex \(2\times2\) matrices, and its actual matrix exponential. The probability coordinates are the diagonal entries in \eqref{eq:readout}; the closed region is the closure of these finite outputs. The correspondence with all four entries of \eqref{eq:stochastic} is proved, so the two-coordinate representation does not replace the original stochastic-table requirement by a weaker condition.

The development also proves the bridges used in the analysis: removing scalar jump parts with the Hamiltonian correction; changing from the complex matrix representation to the real augmented Bloch equation; realizing every positive-semidefinite canonical matrix by finitely many Pauli jumps; and equating the resulting readouts with the original probabilities. The reference first-order integrals are connected to the same exponential derivative used in the global estimates and the analytic system. Thus the spectral calculation, the finite-attainment argument, and the implicit-function theorem apply to the original \(\Q\) and \(f\).

The finite-dimensional channel foundations use a vendored, attributed copy of QICLean\footnote{Source revision: \url{https://github.com/LionSR/QICLean/tree/4d5bcdbe6e92fa64c60d0361cf5472445deba534}. Its Apache-2.0 license and compatibility patch are included in the source distribution.}. The main declared theorems are audited with \texttt{\#print axioms}. Their dependencies are the standard Lean principles \texttt{propext}, \texttt{Classical.choice}, and \texttt{Quot.sound}; no \texttt{sorryAx} or additional mathematical axiom is used. These are exact real and complex proofs: in particular \(\pi\) and \(\sqrt2\) in the swap slope are not replaced by rational numerical surrogates.

\subsection{Statement correspondence}
Table~\ref{tab:formal-map} gives the principal endpoints of the formal proof chain. Except for the namespaces \texttt{Swap}, \texttt{TwoState.GeneralGeometry}, and \texttt{TwoStateRoots}, all declaration names in the table have the prefix \texttt{TwoStateQuantumEmbedding}. The source distribution includes a fuller statement map with the hypotheses and the supporting bridge declarations. Intermediate statements about neighborhoods or local minima retain their hypotheses; the final boundary theorem discharges them using actual minimizers.

\small
\begin{longtable}{@{}>{\raggedright\arraybackslash}p{0.31\textwidth}>{\raggedright\arraybackslash}p{0.65\textwidth}@{}}
\caption{Correspondence between the paper and the formal development.}\label{tab:formal-map}\\
\toprule
Mathematical statement & Principal Lean declarations\\
\midrule
\endfirsthead
\toprule
Mathematical statement & Principal Lean declarations\\
\midrule
\endhead
\bottomrule
\endfoot
Original table and canonical model &
\path{OriginalProblem.original_Q2_four_entry_iff};
\path{CanonicalBloch.finiteOutput_eq_canonicalOutput};
\path{CanonicalBloch.finiteOutputs_eq_canonicalOutputs}\\[3pt]
Theorem~\ref{thm:drift-ellipsoid} &
\path{TwoState.GeneralGeometry.physicalCanonical_posSemidef_iff_ball}\\[3pt]
Proposition~\ref{prop:fixed-drift} &
\path{TwoState.GeneralGeometry.physicalCanonical_linear_image};
\path{OptimalDrift.explicit_endpoint_output}\\[3pt]
Proposition~\ref{prop:filling} &
\path{CanonicalBias.Q2_antidiagonal_filling};
\path{CanonicalBias.Q2_swap}\\[3pt]
Theorem~\ref{thm:envelope} &
\path{Envelope.Q2_envelope_iff};
\path{Envelope.F_upperSemicontinuousOn};
\path{Envelope.F_classical};
\path{Envelope.relaxed_variational_limit}\\[3pt]
Lemma~\ref{lem:rigidity} &
\path{Rigidity.canonicalOutput_dissipative_rigidity};
\path{Rigidity.finiteOutput_dissipative_rigidity}\\[3pt]
Lemma~\ref{lem:uniform-swap} &
\path{SwapComparison.canonical_near_swap_comparison_all_axes};
\path{SwapComparison.finite_near_swap_comparison}\\[3pt]
Lemma~\ref{lem:spectrum} &
\path{Swap.all_frequency_sharp};
\path{Swap.higher_frequency_gap};
\path{Swap.first_frequency_gap}\\[3pt]
Lemma~\ref{lem:matching} &
\path{SwapMatching.lowerBoundary_small_upper_strict}, with the actual matching-output and intermediate-value lemmas in the same module\\[3pt]
Theorem~\ref{thm:swap-cone} &
\path{SwapCone.Q2_quadratic_outer};
\path{SwapCone.Q2_posTangentConeAt_zero}\\[3pt]
Lemma~\ref{lem:localization} &
\path{NearSwap.canonical_near_localization};
\path{NearSwap.canonical_normalized_parameters_tendsto}\\[3pt]
Lemma~\ref{lem:attainment} &
\path{SwapFiniteAttainment.exists_swapMinimizer};
\path{SwapFiniteAttainment.all_minimizers_localized}\\[3pt]
Lemma~\ref{lem:rank-one} and the full chart &
\path{SwapRankOne.eventually_rank_one_of_tendsto};
\path{GlobalSwapArc.minimizing_rank_one_eventually};
\path{GlobalSwapArc.minimizing_parameters_tendsto}\\[3pt]
Actual stationary equations for minimizers &
\path{ActualSwapOptimality.stationarity_chart_of_lowerBoundary};
\path{GlobalSwapArc.minimizing_stationarity_eventually}\\[3pt]
Lemma~\ref{lem:analytic-system} &
\path{ActualSwapSystem.actual_analytic_stationary_branch};
\path{ActualSwapSystem.actual_stationary_arc_reparameterized}\\[3pt]
Theorem~\ref{thm:analytic-boundary} &
\path{GlobalSwapArc.lowerBoundary_analytic_near_swap}\\[3pt]
Proposition~\ref{prop:planarity} &
\path{ActualSwapSystem.unique_stationary_branch_eventually_planar};
\path{ActualSwapSystem.actual_planar_analytic_stationary_branch}, combined with the preceding identification theorem\\[3pt]
Theorem~\ref{thm:reset} &
\path{ResetFamily.reset_every_small_epsilon_uniform}; the actual jump and readout are defined in \path{ResetFamily}\\[3pt]
Theorem~\ref{thm:square-roots} &
\path{TwoStateRoots.rootLayer_two_eq_square}\\
\end{longtable}
\normalsize

For precision, the final analytic-boundary declaration has the mathematical type
\[
 \exists f_{\mathrm{arc}}:\R\to\R,\ \exists b_0>0:
 \quad f_{\mathrm{arc}}\text{ is analytic on }(-b_0,b_0),\qquad
 \forall b\in[0,b_0),\ f_{\mathrm{arc}}(b)=f(b).
\]
There is no assumed local boundary identity in this statement. The formal proof obtains it from finite minimizers, their uniform parameter limits, the actual stationary equations, and the scalar inverse of the actual second probability. Planarity is proved for the unique stationary branch and transferred to its positive-scale boundary realizations. The reset theorem likewise quantifies over every sufficiently small prescribed \(\epsilon\), with one uniform remainder constant; it is not only a limit along selected parameters.

\Needspace{9\baselineskip}
\subsection{Reproduction and remaining geometry}
From the root of the distributed source package, the verification commands are
\begin{verbatim}
cd lean
lake exe cache get
lake build TwoStateQuantumEmbedding QICLean
lake env lean ../verification/PublicationAudit.lean
\end{verbatim}
The toolchain and mathlib revision are pinned in the project. The audit file names the publication endpoints and prints their axiom dependencies. The paper sources are self-contained LaTeX files and do not depend on the research workspace.

The results determine the exact first-order geometry at swap and prove an actual analytic boundary on a positive neighborhood of that endpoint. They also give global filled sections and the classical half-region. They do not determine the entire nonclassical envelope for \(-1<x<0\). In particular, this paper makes no claim of global continuity of \(F\), finite generator attainment at every negative \(x\), a universal reset upper bound with coefficient \(4/e\), or a classification of higher channel-root layers. The reset result is a constructive lower improvement, and \(\mathcal R_2=[0,1]^2\) is a statement about squares alone. These distinctions keep the established local boundary theorem separate from the remaining global embedding problem.

\begin{thebibliography}{9}
\bibitem{GKS1976}
V.~Gorini, A.~Kossakowski, and E.~C.~G.~Sudarshan,
Completely positive dynamical semigroups of \(N\)-level systems,
\emph{Journal of Mathematical Physics} \textbf{17}, 821--825 (1976).
\href{https://doi.org/10.1063/1.522979}{doi:10.1063/1.522979}.

\bibitem{Lindblad1976}
G.~Lindblad,
On the generators of quantum dynamical semigroups,
\emph{Communications in Mathematical Physics} \textbf{48}, 119--130 (1976).
\href{https://doi.org/10.1007/BF01608499}{doi:10.1007/BF01608499}.

\bibitem{Shahbeigi2024}
F.~Shahbeigi, C.~T.~Chubb, R.~Kukulski, {\L}.~Pawela, and K.~Korzekwa,
Quantum-embeddable stochastic matrices,
\emph{Quantum} \textbf{8}, 1404 (2024).
\href{https://doi.org/10.22331/q-2024-07-10-1404}{doi:10.22331/q-2024-07-10-1404}.

\bibitem{Korzekwa2018}
K.~Korzekwa, S.~Czach\'orski, Z.~Pucha\l a, and K.~\.{Z}yczkowski,
Coherifying quantum channels,
\emph{New Journal of Physics} \textbf{20}, 043028 (2018).
\href{https://doi.org/10.1088/1367-2630/aaaff3}{doi:10.1088/1367-2630/aaaff3}.

\bibitem{WolfCirac2008}
M.~M.~Wolf and J.~I.~Cirac,
Dividing quantum channels,
\emph{Communications in Mathematical Physics} \textbf{279}, 147--168 (2008).
\href{https://doi.org/10.1007/s00220-008-0411-y}{doi:10.1007/s00220-008-0411-y}.
\end{thebibliography}
\end{document}



